\section{Cierre: un solo valor, un solo juego, y la entrega como teorema}
\label{sec:cierre}

Las partes anteriores construyen piezas —potencial híbrido, mercado de
\emph{wrench}, juego de trayectorias, testigo, barreras, arrendamiento— y cada
una tiene su garantía. Falta el enunciado que las une: que el mecanismo
distribuido, ejecutado por robots con ruedas y contacto e información de $n$
saltos, entrega la carga. Esta sección lo enuncia con sus hipótesis, lo
demuestra por composición, y en \cref{sec:ensayo-cierre} lo ensaya sobre la
planta con el mecanismo distribuido y no con un controlador central.

\subsection{El potencial es el valor de misión}

El potencial de \cref{th:potencial} y el valor $V_H^\star$ de
\cref{def:continuaciones} se enunciaron como objetos distintos. No lo son.

\begin{definicion}[Valor de misión unificado]
\label{def:valor-unificado}
Sea $z=(d,c)$ el perfil híbrido: parte discreta $d=(d_i)_i$ con
$d_i\in\{\mathrm{idle}\}\cup\{(k,\mu,s)\}$ y parte continua $c$ (esfuerzos,
continuaciones). Para cada carga $k$ sea $\Ccal_k(d)$ la coalición que $d$ le
asigna, $\mathcal T_k(\Ccal_k,Y)$ sus continuaciones ejecutables
\eqref{eq:T-continuaciones} y $V_k(\Ccal_k,Y)=\sup_{\tau\in\mathcal T_k}
\bigl[B_k\,\mathbf 1_{\text{entrega}}(\tau)-\mathcal H_k(\tau)\bigr]$ su valor
(con $V_k=-\infty$ si $\mathcal T_k=\emptyset$ y $V_k$ definido por
prospecto creciente en $|\Ccal_k|$ cuando $|\Ccal_k|<r_k^\star$). El valor de
misión es
\begin{equation}
  \label{eq:valor-unificado}
  \Vcal(z,Y)=\sum_k V_k(\Ccal_k(d),Y)-\sum_i c_i(d_i,Y)-\Gamma(z,Y),
\end{equation}
con $c_i$ el coste privado de $i$ (viaje, batería) y $\Gamma$ los
acoplamientos simétricos (congestión atómica de
\cref{th:externalidad-atomica}, exclusividad de puesto, arrendamientos).
\end{definicion}

\begin{teorema}[El potencial híbrido es el valor de misión]
\label{th:potencial-valor}
Si el pago de cada AMR se construye como utilidad marginal del valor
unificado, $J_i(z)=\Vcal(z,Y)-\Vcal(z_i^{0},z_{-i},Y)$ con
$z_i^{0}=\mathrm{idle}$, entonces $\Phi\equiv\Vcal$ es un potencial exacto del
megajuego en el sentido de \cref{th:potencial}: para toda desviación
unilateral factible, discreta, continua o mixta,
$J_i(z_i',z_{-i})-J_i(z_i,z_{-i})=\Vcal(z_i',z_{-i})-\Vcal(z_i,z_{-i})$. En
particular, dentro de una rama discreta fija $d$, el juego continuo tiene
potencial $\Vcal(d,\cdot)$, que coincide con $-$ la suma de valores
$V_H^\star$ de \cref{def:continuaciones} salvo la constante $\sum_kB_k$; y el
equilibrio de \cref{pr:commit} y el de \cref{th:externalidad-atomica} son el
mismo objeto medido en la misma unidad.
\end{teorema}

\begin{proof}
El término $\Vcal(z_i^0,z_{-i},Y)$ no depende de $z_i$ y se cancela en la
diferencia, exactamente como en \cref{th:potencial}; que $\Vcal$ sea el
potencial es entonces la definición. La identificación con $V_H^\star$ es
\eqref{eq:valor-mision} con $\mathcal H=\mathcal H_k$ y la indicatriz de
entrega, sumada sobre cargas. La externalidad atómica de $\Gamma$ es
simétrica por \cref{th:externalidad-atomica}, luego no rompe la exactitud.
\end{proof}

\begin{observacion}[Lo que esto cierra y lo que no]
\label{obs:potencial-valor-alcance}
Cierra la objeción de que reclutamiento, contacto, ruta y reemplazo se
comparasen en unidades distintas: \cref{obs:coherencia-valor} deja de ser un
principio y pasa a ser una identidad. No cierra la optimalidad: el máximo de
$\Vcal$ es el óptimo de misión, pero el mecanismo solo garantiza llegar a un
equilibrio de orden $h$ de $\Vcal$, y \cref{th:descomposicion} mide la
distancia.
\end{observacion}

\subsection{El teorema de entrega}

\begin{supuesto}[Hipótesis del cierre]
\label{s:cierre}
\begin{enumerate}[label=(A\arabic*),leftmargin=2.6em,topsep=0.2em,itemsep=0.1em]
  \item \emph{Finitud y valor.} El conjunto de partes discretas es finito y
  $\Vcal$ es el valor unificado de \cref{def:valor-unificado}, con revisiones
  aceptadas solo si mejoran $\Vcal$ al menos $\epsilon_{\mathrm{sw}}>0$
  (histéresis de \cref{th:dwell}), certificadas por \cref{pr:commit}.
  \item \emph{Certificado físico.} Una coalición pasa a formación solo si su
  margen de \emph{wrench} es no negativo con la geometría de puestos declarada
  (\cref{th:span}, \cref{th:cert}) y, en \textsf{caging}, si la formación
  cierra (\cref{th:cerco}); tras la formación se recertifica con las
  capacidades observadas (\cref{obs:certificado-postevento}).
  \item \emph{Mercado y rama convexa.} Dentro de cada rama, el reparto de
  esfuerzos es el reposo de \cref{th:N-precios} con paso admisible
  (\cref{th:N-ventana}) y los límites de contacto son convexos (disco en
  \textsf{cargo}, semieje en \textsf{caging}).
  \item \emph{Testigo y progreso.} Cada coalición en transporte mantiene una
  continuación incumbente cuyo presupuesto $W$ decrece al menos $\rho_0$ por
  unidad de tiempo físico fuera de las paradas seguras
  (\cref{pr:presupuesto-progreso}), y el corredor se arrienda con orden total
  (\cref{pr:arrendamiento}, \cref{th:ciclo-espera}).
  \item \emph{Seguridad.} Las barreras de coalición y de AMR son factibles en
  todo estado alcanzado con el estrechamiento de \cref{le:muestras}, y una
  parada segura existe desde cualquier estado alcanzado.
\end{enumerate}
\end{supuesto}

\begin{teorema}[Entrega segura por el mecanismo distribuido]
\label{th:entrega}
Bajo \cref{s:cierre}, con información de $n$ saltos y sin ningún nodo que
conozca el estado completo:
\begin{enumerate}[label=(\roman*),leftmargin=2.2em,topsep=0.2em,itemsep=0.1em]
  \item la capa discreta termina: el número de revisiones aceptadas en toda la
  misión es a lo sumo $(\Vcal_{\max}-\Vcal(z_0))/\epsilon_{\mathrm{sw}}$, y
  entre dos revisiones de un mismo AMR media al menos
  $2\epsilon_{\mathrm{sw}}/L_{\mathrm{sw}}$;
  \item toda coalición que pasa a formación tiene margen de \emph{wrench} no
  negativo con las capacidades reales, luego existe un reparto de esfuerzos
  admisible para la demanda de diseño, y el mercado converge a él;
  \item ninguna barrera se viola y no hay interbloqueo en el corredor;
  \item cada carga cuya coalición entra en transporte con presupuesto $W_0$ es
  entregada antes de $W_0/\rho_0$ unidades de tiempo de avance, y la misión
  termina en tiempo finito con todas las cargas certificables entregadas.
\end{enumerate}
\end{teorema}

\begin{proof}
(i) Por \cref{th:potencial-valor}, cada revisión aceptada aumenta $\Vcal$ al
menos $\epsilon_{\mathrm{sw}}$ (\cref{pr:commit} garantiza que la mejora
certificada es real); $\Vcal$ está acotado superiormente sobre un conjunto
finito de partes discretas con óptimo continuo único por rama (A3), luego el
número de aceptaciones es finito con la cota dada, y la permanencia mínima es
\cref{th:dwell}. (ii) La certificación de (A2) con capacidades observadas es
\cref{th:span}: margen no negativo equivale a que la demanda de diseño
pertenezca al conjunto realizable, que es convexo por (A3); el reposo de
\cref{th:N-precios} existe, es único y la iteración con el paso de
\cref{th:N-ventana} converge a él; en \textsf{caging} la unilateralidad entra
en el conjunto realizable como semiejes y \cref{th:cerco} garantiza que la
carga no escapa de la formación. (iii) Con (A5), \cref{pr:hocbf} y
\cref{le:muestras} dan invariancia del conjunto seguro bajo ejecución
muestreada; el arrendamiento con orden total excluye el ciclo de espera por
\cref{th:ciclo-espera}, luego no hay interbloqueo, y \cref{pr:arrendamiento}
excluye dos coaliciones en el corredor. (iv) Fuera de las paradas seguras el
presupuesto decrece al menos $\rho_0$ por unidad de tiempo
(\cref{pr:presupuesto-progreso}); las paradas seguras son finitas en número y
duración porque solo las produce la espera del arrendamiento, que termina por
(iii). Al agotarse $W$ la carga está en el conjunto terminal por definición
del testigo (\cref{th:testigo}); las coaliciones liberadas vuelven al
reclutamiento y (i) acota las revisiones restantes.
\end{proof}

\begin{observacion}[Qué no dice el teorema, y por qué es útil de todos modos]
\label{obs:entrega-alcance}
No dice que la misión sea óptima: dice que termina, con entrega, sin
colisión y sin interbloqueo, bajo cinco hipótesis que son cada una un
enunciado anterior del documento. Tampoco dice que (A5) sea gratis: la
existencia de una parada segura desde todo estado alcanzado es la hipótesis
fuerte, y \cref{th:grado-relativo-rueda} dice que la barrera correcta para la
carga es de cuarto orden en el par. Lo que el teorema aporta es la lista de
lo que hay que verificar en una planta concreta para reclamar la entrega, y
esa lista es la que \cref{sec:ensayo-cierre} recorre.
\end{observacion}
